\documentclass[12pt]{article}

\usepackage[utf8]{inputenc}
\usepackage{amsmath, amssymb, amsthm}
\usepackage{geometry}
\geometry{a4paper, margin=2.5cm}

\title{Generatori și Structuri Generate}
\author{}
\date{}

\theoremstyle{definition}
\newtheorem{definition}{Definiție}

\begin{document}

\maketitle

\section{Generatori}

\begin{definition}
Fie $\Omega$ un spațiu și $\mathcal{G} \subseteq \mathcal{P}(\Omega)$ o familie de mulțimi.
Spunem că $\mathcal{G}$ este un \emph{generator} pentru o anumită clasă de structuri (algebre, sigma-algebre, lambda-sisteme, clase monotone, $\pi$-sisteme) dacă structura dorită se obține ca cea mai mică structură de acel tip care conține $\mathcal{G}$, adică drept intersecția tuturor structurilor de acel tip care conțin $\mathcal{G}$.
\end{definition}

\noindent
Formal, dacă $\mathcal{C}$ este o clasă de structuri (de exemplu toate sigma-algebrele de pe $\Omega$), atunci structura generată de $\mathcal{G}$ este


\[
    \bigcap \{\mathcal{S} \in \mathcal{C} : \mathcal{G} \subseteq \mathcal{S}\}.
\]



\section{Algebra generată}

\begin{definition}
Fie $\mathcal{G} \subseteq \mathcal{P}(\Omega)$. \emph{Algebra generată} de $\mathcal{G}$ este


\[
    \operatorname{Alg}(\mathcal{G})
    :=
    \bigcap \{\mathcal{A} : \mathcal{A} \text{ este algebră de mulțimi și } \mathcal{G} \subseteq \mathcal{A}\}.
\]


Este cea mai mică algebră de mulțimi care conține familia $\mathcal{G}$.
\end{definition}

\section{Sigma-algebra generată}

\begin{definition}
Fie $\mathcal{G} \subseteq \mathcal{P}(\Omega)$. \emph{Sigma-algebra generată} de $\mathcal{G}$ este


\[
    \sigma(\mathcal{G})
    :=
    \bigcap \{\mathcal{F} : \mathcal{F} \text{ este sigma-algebră și } \mathcal{G} \subseteq \mathcal{F}\}.
\]


Este cea mai mică sigma-algebră care conține familia $\mathcal{G}$.
\end{definition}

\section{Lambda-sistemul generat}

\begin{definition}
Fie $\mathcal{G} \subseteq \mathcal{P}(\Omega)$. \emph{Lambda-sistemul generat} de $\mathcal{G}$ este


\[
    \lambda(\mathcal{G})
    :=
    \bigcap \{\mathcal{D} : \mathcal{D} \text{ este lambda-sistem și } \mathcal{G} \subseteq \mathcal{D}\}.
\]


Este cel mai mic lambda-sistem care conține familia $\mathcal{G}$.
\end{definition}

\section{Clasa monotonă generată}

\begin{definition}
Fie $\mathcal{G} \subseteq \mathcal{P}(\Omega)$. \emph{Clasa monotonă generată} de $\mathcal{G}$ este


\[
    \mathcal{M}(\mathcal{G})
    :=
    \bigcap \{\mathcal{M} : \mathcal{M} \text{ este clasă monotonă și } \mathcal{G} \subseteq \mathcal{M}\}.
\]


Este cea mai mică clasă monotonă care conține familia $\mathcal{G}$.
\end{definition}

\section{Pi-sistemul generat}

\begin{definition}
Fie $\mathcal{G} \subseteq \mathcal{P}(\Omega)$. \emph{Pi-sistemul generat} de $\mathcal{G}$ este


\[
    \pi(\mathcal{G})
    :=
    \bigcap \{\mathcal{P} : \mathcal{P} \text{ este pi-sistem și } \mathcal{G} \subseteq \mathcal{P}\}.
\]


Este cel mai mic pi-sistem (adică familie închisă la intersecții finite) care conține familia $\mathcal{G}$.
\end{definition}

\section{Relații între structurile generate}

\noindent
Dacă $\mathcal{G}$ este un $\pi$-sistem, atunci sunt valabile relațiile clasice:
\begin{itemize}
    \item $\sigma(\mathcal{G}) = \lambda(\mathcal{G})$ (teorema Dynkin--$\pi$),
    \item $\sigma(\mathcal{G}) = \mathcal{M}(\mathcal{G})$ (teorema clasei monotone).
\end{itemize}

\end{document}
